\documentclass[11pt]{article}

\usepackage[margin=1in]{geometry}
\usepackage{amsmath}
\usepackage{amssymb}
\usepackage{amsthm}
\usepackage{graphicx}
\usepackage{booktabs}
\usepackage{siunitx}
\usepackage{microtype}
\usepackage{natbib}
\usepackage{hyperref}
\usepackage{cleveref}

\title{Working title}
\author{Author One\\ Institution\\ \texttt{author@example.org}}
\date{}

\begin{document}

\maketitle

\begin{abstract}
One paragraph. State the problem, what you did, the single most important
result with a number, and why it matters. Write this last.
\end{abstract}

\section{Introduction}
\label{sec:introduction}

What question this paper answers, why it is open, and what the paper
contributes. End with an explicit list of contributions.

\section{Related work}
\label{sec:related}

The lines of work this paper sits between, and what each one leaves unsolved.

\section{Method}
\label{sec:method}

The setup, the notation, and the procedure. Define every symbol before using it.

\begin{equation}
\label{eq:objective}
\mathcal{L}(\theta) = \frac{1}{n} \sum_{i=1}^{n} \ell\bigl(f_\theta(x_i), y_i\bigr).
\end{equation}

\Cref{eq:objective} is the objective optimized throughout.

\section{Experiments}
\label{sec:experiments}

Data, baselines, protocol, and metrics, in that order. Report uncertainty and
say what it is.

\begin{table}[t]
  \centering
  \caption{Results on the evaluation suite. Replace with real numbers.}
  \label{tab:results}
  \begin{tabular}{lcc}
    \toprule
    Method & Accuracy (\%) & Latency (ms) \\
    \midrule
    Baseline & 00.0 & 000 \\
    This work & 00.0 & 000 \\
    \bottomrule
  \end{tabular}
\end{table}

\section{Conclusion}
\label{sec:conclusion}

What is now known that was not known before, and the concrete limitation a
reader should carry away.

\bibliographystyle{plainnat}
\bibliography{references}

\end{document}
